\section{The sharp triangular clock}
View a composition as consecutive intervals partitioning $N$ labelled unit
cells. Mass means interval length. Every update coarsens this interval
partition, so a deleted cut never returns. A block is \emph{born in round
$t$} if it combines at least two blocks when time $t-1$ becomes time $t$;
an unchanged block retains its earlier identity.

The fixed points are exactly the strictly decreasing compositions. Indeed,
all runs are singletons precisely in that case. Otherwise an update
strictly decreases the number of parts, proving termination. To sharpen
this generic bound, we first control the left side of a delayed cut.

\begin{lemma}[Left mass and right birth]\label{lem:leftmass}
If a cut disappears in round $t\ge1$, its old left block has mass at least
$t$. If $t\ge2$, its old right block was born in round $t-1$.
\end{lemma}
\begin{proof}
Positivity gives the mass bound for $t=1$. For $t\ge2$, let $A,B$ be the
adjacent time-$(t-1)$ blocks at the disappearing cut, so
$\wt A\le\wt B$. At time $t-2$, let $a$ be the rightmost parent of $A$
and $c$ the leftmost parent of $B$. Their boundary survived round $t-1$,
hence $\wt a>\wt c$. If $B$ had not been born then, it would equal $c$,
and $\wt A\ge\wt a>\wt c=\wt B$, a contradiction. Thus $B$ was born in
round $t-1$. Its first internal cut had left block $c$ and disappeared in
that round. Induction gives $\wt c\ge t-1$; integrality now yields
$\wt A\ge\wt a\ge\wt c+1\ge t$. This proves both assertions.
\end{proof}

\begin{lemma}[New-block mass]\label{lem:birthmass}
Every block born in round $t\ge1$ has mass at least
$M_t=1+t(t+1)/2$.
\end{lemma}
\begin{proof}
A block born in round $1$ has at least two positive parents, so its mass is
at least $2=M_1$. For $t\ge2$, take the first two parents $A_1,A_2$ of a
new block. Their cut disappears in round $t$. Lemma~\ref{lem:leftmass}
gives $\wt{A_1}\ge t$ and says that $A_2$ was born in round $t-1$.
The induction hypothesis gives $\wt{A_2}\ge M_{t-1}$. The parents are
disjoint, so the new mass is at least $t+M_{t-1}=M_t$. Additional parents
only increase it.
\end{proof}

\begin{theorem}[Sharp clock]\label{thm:clock}
For every $N\ge1$,
\[
 \max_{a\in\Comp_N}\tau(a)
 = H(N):=\max\{h\ge0:1+h(h+1)/2\le N\}.
\]
More generally, for every $h\ge1$ and $r\ge0$, the composition
$(h,h-1,\ldots,1,1+r)$ has stabilization time exactly $h$.
\end{theorem}
\begin{proof}
If $\tau(a)=t\ge1$, its last nonfixed round creates a block, so
Lemma~\ref{lem:birthmass} gives $N\ge M_t$. This proves the upper bound.
For the witness, after $j$ updates, $1\le j\le h$, the state is
\begin{equation}\label{eq:witness}
 \bigl(h,h-1,\ldots,j+1,\;r+1+j(j+1)/2\bigr),
\end{equation}
where the prefix is empty for $j=h$. At the first update only the final
pair $(1,1+r)$ merges. If \eqref{eq:witness} holds with $j<h$, its suffix
mass is at least $j+1$, and every earlier old prefix pair is a strict
descent. Exactly the last prefix block joins the suffix, giving the formula
for $j+1$. Thus exactly $h$ nonfixed updates lead to a single part. This
argument permits every surplus $r$, however large. For $N\ge2$ choose
$h=H(N)$ and $r=N-M_h$; for $N=1$, the unique state $(1)$ has time zero.
\end{proof}
